\subsection{Random Polygon Algorithm} % (fold)

\begin{figure}[h]
\begin{center}
\includegraphics[width=0.3\textwidth]{img/rpa200.eps}
\end{center}
\caption{Polygon mit 200 Punkten (Random Polygon Algorithm)}
\label{fig:virmani200}
\end{figure}

\label{sub:random_polygon_algorithm}

  Der \textit{Random Polygon Algorithm} (RPA) ist ein von Dailey und
  Whitfield in~\cite{dailey08rpa} beschriebener Algorithmus zur Erzeugung
  beliebiger Repräsentanten aus der Klasse aller einfachen Polygone.

  \subsubsection{Algorithmus} % (fold)
  \label{ssub:algorithmus}

    Der RPA arbeitet im Gegensatz zu den anderen hier aufgefürten
    Algorithmen nicht auf einer gegebenen Punktmenge. Stattdessen werden
    alle Punkte des erzeugten Polygons zur laufzeit berechnet und
    eingefügt. Zur initialisierung werden zunächst drei beliebige Punkte
    innerhalb der Bounding Box gewählt. Anschliessend werden für ein n-gon
    die folgenden Schritte (n-3)-mal ausgeführt:

    \begin{itemize}
      \item Wähle eine zufällige Kante $V_aV_b$ des Polygons.
      \item Bestimme die von dieser Kante sichtbare Region.
      \item Wähle einen zufälligen Punkt $V_c$ innerhalb der 
            erzeugten sichtbaren Region und füge diesen in das 
            Polygon ein. Ersetze die Kante $V_aV_b$ durch die Kanten 
            $V_aV_c$ und $V_cV_b$.
    \end{itemize}

    Zur Bestimmung der sichtbaren Region wird als erstes die von der Kante 
    $V_aV_b$
    sichtbare Region innerhalb des Polygons ermittelt. Anschliessend wird die
    von $V_aV_b$ sichtbare Region ausserhalb des Polygons und innerhalb der
    Bounding Box bestimmt und mit der inneren verbunden. Zur Erzeugung des
    Punktes $V_c$ wird die erhaltene Gesamtregion trianguliert und aus der
    erhaltenen Menge von Dreiecken eins gewichtet nach der Größe zufällig
    gewählt. Innerhalb dieses Dreiecks wird nun ein Punkt zufällig gewält und
    als Punkt $V_c$ in das Polygon eingefügt. Für eine detailierte Beschreibung
    des Algorithmus siehe~\cite{dailey08rpa}.
  
  % subsubsection algorithmus (end)

  \subsubsection{Eigenschaften \& Implementierungsdetails} % (fold)
  \label{ssub:eigenschaften}

    Um den inneren und äußeren Teil der von $V_aV_b$ sichtbaren Region
    aufzubauen ist es nötig zu bestimmen ob ein Punkt von $V_a$ und von $V_b$
    aus sichtbar sind. Zudem werden sichtbare Punkte in von innen sichtbare
    (d.h. die Sichtlinie befindet sich gänzlich innerhalb des Polygons) und
    von aussen sichtbare (d.h. die Sichtlinie befindet sich gänzlich
    ausserhalb des Polygons, jedoch innerhalb der Boundingbox) unterteilt.
    In~\cite{dailey08rpa} wird hierzu ein einfacher Test auf Schnittpunkte mit
    der Kante $V_aV_b$ vorgeschlagen, jedoch nicht beschrieben, mit welcher
    weiteren Kante $V_aV_b$ geschnitten wird. Für einen Punkt $V_x$ ist es
    jedoch mit einem Sichtbarkeitstest und zwei Winkeltest möglich zu
    bestimmen, ob der Weg auf dem Polygon von $V_a$ bzw. $V_b$ zu $V_x$ links
    oder rechts der Sichtlinie von $V_a$ bzw. $V_b$ zu $V_x$ liegt.

    Die von Dailey und Whitfield in~\cite{dailey08rpa} abgeschätzte worst-case
    Laufzeit des RPA beträgt $\bigO(n^3)$. Hierbei gehen sie von einer
    Triangulierung in $\bigO(n \log n)$ aus. Der von uns verwendete Ear
    Clipping Algorithmus erreicht jedoch nur eine Laufzeit von $\bigO(n^2)$,
    so dass die worst-case Laufzeit des RPA insgesamt auf $\bigO(n^4)$ steigt.
    Mit einem schnelleren, jedoch gleichzeitig wesentlich komplexeren
    Triangulierungsalgorithmus wie dem von Seidel~\cite{seidel91asimple} wäre
    die von Dailey und Whitfiel abgeschätzte Laufzeit auch mit der
    vorliegenden Implementierung des RPA zu erreichen.

    Von den von uns implementierten Algorithmen ist der RPA der einzige,
    welcher Repräsentanten aus der Klasse aller einfachen Polygone erzeugt.
    Der von Dailey und Whitfield hierzu geführete Beweis ist in
    ~\cite{dailey08rpa} zu finden.

  % subsubsection eigenschaften (end)

% subsection random_polygon_algorithm (end)